\origpage{363}{N. Robertson, P.D. Seymour / Journal of Combinatorial Theory, Series B 90 (2004) 325--385}
From (1), it follows that if $k \le \frac{1}{2}\rho$ and $r_1, \dots, r_k$ is a sequence of regions of $\Gamma$ such that $\bar{r}_1 \cap U(H) \ne \emptyset$ and $\bar{r}_i \cap \bar{r}_{i+1} \ne \emptyset$ for $1 \le i < k$ then $\bar{r}_k \subseteq \Lambda$. Consequently the same statement holds for $sk(\Gamma)$. Since $\rho \ge 2$, and $U(H)$ is non-empty and arc-wise connected, it follows from Theorem 5.2 of [8] that there is a circuit $C_1$ of $sk(\Gamma)$ with $U(C_1) \subseteq \Lambda$, bounding an open disc in $\Lambda$ including $U(H)$, such that every edge of $C_1$ is incident with a region $r$ of $sk(\Gamma)$ satisfying $\bar{r} \cap U(H) \ne \emptyset$. By (1) and $\frac{1}{2}\rho - 1$ further applications of Theorem 5.2 of [8], there are circuits $C_1, \dots, C_{\frac{1}{2}\rho}$ of $sk(\Gamma)$, such that for $2 \le i \le \frac{1}{2}\rho$, $U(C_i)$ bounds an open disc in $\Lambda$ including $U(C_{i-1})$, and every edge of $C_i$ is incident with a region $r$ of $sk(\Gamma)$ satisfying $\bar{r} \cap U(C_{i-1}) \ne \emptyset$. We deduce

\hspace*{12pt}(2) \textit{For each $v \in V(C_{\frac{1}{2}\rho})$ there is a path of $\Gamma^*$ from $U(H)$ to $v$ with $\le \rho$ edges, and every such path has $\ge \rho - 1$ edges.}

\textbf{Subproof.} Let $v \in V(C_{\frac{1}{2}\rho})$. We have seen, from the construction of $C_1, \dots, C_{\frac{1}{2}\rho}$, that there is a sequence $r_1, \dots, r_{\frac{1}{2}\rho}$ of regions of $sk(\Gamma)$ such that $\bar{r}_1 \cap U(H) \ne \emptyset$, $v \in \bar{r}_{\frac{1}{2}\rho}$, and $\bar{r}_i \cap \bar{r}_{i+1} \ne \emptyset$ for $1 \le i < \frac{1}{2}\rho$. But for any two vertices of $\Gamma^*$, if there is a region of $sk(\Gamma)$ whose closure contains them both, then there is also a region of $\Gamma$ whose closure contains them both. Hence we may assume that $r_1, \dots, r_{\frac{1}{2}\rho}$ are regions of $\Gamma$; and so there is a path of $\Gamma^*$ from $U(H)$ to $v$ with $\le \rho$ edges. Any such path $P$ meets $V(C_1), \dots, V(C_{\frac{1}{2}\rho})$, and $|V(P) \cap V(\Gamma)| \le \frac{1}{2}(|E(P)| + 1)$, and so $\frac{1}{2}(|E(P)| + 1) \ge \frac{1}{2}\rho$, that is, $|E(P)| \ge \rho - 1$. This proves (2).

Let $F$ be a circuit of $\Gamma^*$ with $U(F) \subseteq \bar{\Lambda}$, bounding a disc $D(F)$ in $\bar{\Lambda}$ including $U(C_{\frac{1}{2}\rho})$, such that $U(F) \cap V(\Gamma) \subseteq V(C_{\frac{1}{2}\rho})$; and subject to that with $D(F)$ maximal. From the maximality of $D(F)$ it follows that $D(F)$ contains every vertex of $\Gamma^*$ with two neighbours in $V(F) \cap V(C_{\frac{1}{2}\rho})$, and $D(F)$ includes every edge of $\Gamma^*$ with both ends in $V(F)$. It follows that $D(F)$ is a wheel and $H$ is a $D(F)$-hub of radius $\rho$. This proves 12.1. $\square$

We recall that $(\Gamma, \phi, \lambda, \mathscr{T}) \in \mathscr{D}$ is $(m, n, \omega_0)$-flawed internally. Let $Z = \{e \in E(\Gamma) : \phi(e) \ge \omega_0\}$. Since $\chi$ is disjoint, every member of $Z$ is internal. We denote the set of all poles of $\Gamma^*$ by $P^*$.

\textbf{12.2.} \textit{There exists $X_1 \subseteq V(\Gamma^*)$ with $|X_1| \le m + c(\Sigma_\chi)$ and $P^* \subseteq X_1$, such that for every $e \in Z$ there exists $x \in X_1$ with $d(e, x) \le n$.}

\textbf{Proof.} Choose $e_1, \dots, e_k \in Z$ with $k$ maximum such that $d(e_i, e_j) > n$ ($1 \le i < j \le k$) and $d(e_i, r^*) > n$ ($1 \le i \le k$) for every pole $r^*$. Since $(\Gamma, \phi, \lambda, \mathscr{T})$ is $(m, n, \omega_0)$-flawed internally, it follows that $k < m$. For $1 \le i \le k$ let $x_i \in \tilde{e}_i$, and let $X_1 = P^* \cup \{x_1, \dots, x_k\}$. The result follows. $\square$
